\chapter{2025年全国II卷}

\section{单选题}

\begin{problem}
样本数据 \(2,8,14,16,20\) 的平均数为（\quad）

\choices
    {\(8\)}
    {\(9\)}
    {\(12\)}
    {\(18\)}

\end{problem}

\begin{answer}
C
\end{answer}

\begin{solution}
样本数据 \(2,8,14,16,20\) 的平均数为
\(\frac{2+8+14+16+20}{5}=12\).
故选：C.
\end{solution}

\begin{problem}
已知 \(z=1+\i\)，则 \(\frac{1}{z-1}=\)（\quad）

\choices
    {\(-\i\)}
    {\(\i\)}
    {\(-1\)}
    {\(1\)}

\end{problem}

\begin{answer}
A
\end{answer}

\begin{solution}
因为 \(z=1+\i\)，所以 \(z-1=\i\)，故
\(\frac{1}{z-1}=\frac{1}{\i}=-\i\).
故选：A.
\end{solution}

\begin{problem}
已知集合 \(A=\{-4,0,1,2,8\}\)，\(B=\{x \mid x^3=x\}\)，则 \(A\cap B=\)（\quad）

\choices
    {\(\{0,1,2\}\)}
    {\(\{1,2,8\}\)}
    {\(\{2,8\}\)}
    {\(\{0,1\}\)}

\end{problem}

\begin{answer}
D
\end{answer}

\begin{solution}
由 \(x^3=x\) 得
\(x(x-1)(x+1)=0\)，
所以 \(B=\{-1,0,1\}\). 因此
\(A\cap B=\{0,1\}\).
故选：D.
\end{solution}

\begin{problem}
不等式 \(\frac{x-4}{x-1}\ge 2\) 的解集是（\quad）

\choices
    {\(\{x \mid -2\le x\le 1\}\)}
    {\(\{x \mid x\le -2\}\)}
    {\(\{x \mid -2\le x<1\}\)}
    {\(\{x \mid x>1\}\)}

\end{problem}

\begin{answer}
C
\end{answer}

\begin{solution}
原不等式等价于
\(\frac{x-4}{x-1}-2\ge 0\)，
即
\(\frac{-x-2}{x-1}\ge 0\).
解得 \(-2\le x<1\). 故选：C.
\end{solution}

\begin{problem}
在 \(\triangle ABC\) 中，\(BC=2\)，\(AC=1+\sqrt3\)，\(AB=\sqrt6\)，则 \(A=\)（\quad）

\choices
    {\(45^\circ\)}
    {\(60^\circ\)}
    {\(120^\circ\)}
    {\(135^\circ\)}

\end{problem}

\begin{answer}
A
\end{answer}

\begin{solution}
由余弦定理，
\[
\cos A=\frac{AB^2+AC^2-BC^2}{2AB\cdot AC}
=\frac{6+\left(1+\sqrt3\right)^2-4}{2\sqrt6\left(1+\sqrt3\right)}
=\frac{\sqrt2}{2}
\]
又 \(0^\circ<A<180^\circ\)，故 \(A=45^\circ\). 故选：A.
\end{solution}

\begin{problem}
设抛物线 \(C:y^2=2px\)（\(p>0\)）的焦点为 \(F\)，点 \(A\) 在 \(C\) 上，过 \(A\) 作 \(C\) 的准线的垂线，垂足为 \(B\). 若直线 \(BF\) 的方程为 \(y=-2x+2\)，则 \(|AF|=\)（\quad）

\choices
    {\(3\)}
    {\(4\)}
    {\(5\)}
    {\(6\)}

\end{problem}

\begin{answer}
C
\end{answer}

\begin{solution}
由直线 \(BF:y=-2x+2\)，令 \(y=0\)，得 \(x=1\)，所以 \(F(1,0)\). 故 \(p=2\)，抛物线方程为
\(y^2=4x\)，
其准线方程为 \(x=-1\).

\bitmapfigure[max width=3.8cm,max totalheight=4.8cm]{img/2025/national_paper_2/q6_solution_fig1.png}

因为点 \(B\) 在准线 \(x=-1\) 上，又点 \(B\) 在直线 \(BF:y=-2x+2\) 上，所以
\(B(-1,4)\).
故点 \(A\) 的纵坐标为 \(4\). 代入 \(y^2=4x\)，得
\(x_A=4\).
所以 \(A(4,4)\).

于是
\[
\lvert AF\rvert=\lvert AB\rvert=x_A+\frac p2=4+1=5
\]
故选：C.
\end{solution}

\begin{problem}
记 \(S_n\) 为等差数列 \(\{a_n\}\) 的前 \(n\) 项和，若 \(S_3=6\)，\(S_5=-5\)，则 \(S_6=\)（\quad）

\choices
    {\(-20\)}
    {\(-15\)}
    {\(-10\)}
    {\(-5\)}

\end{problem}

\begin{answer}
B
\end{answer}

\begin{solution}
设等差数列 \(\{a_n\}\) 的公差为 \(d\)，则 \(3a_1+3d=6\)，\(5a_1+10d=-5\).
解得 \(a_1=5\)，\(d=-3\). 所以
\[
S_6=6a_1+15d=6\times 5+15\times(-3)=-15
\]
故选：B.
\end{solution}

\begin{problem}
已知 \(0<\alpha<\pi\)，\(\cos \frac{\alpha}{2}=\frac{\sqrt5}{5}\)，则 \(\sin\left(\alpha-\frac{\pi}{4}\right)=\)（\quad）

\choices
    {\(\frac{\sqrt2}{10}\)}
    {\(\frac{\sqrt2}{5}\)}
    {\(\frac{3\sqrt2}{10}\)}
    {\(\frac{7\sqrt2}{10}\)}

\end{problem}

\begin{answer}
D
\end{answer}

\begin{solution}
由二倍角公式，
\[
\cos\alpha=2\cos^2\frac{\alpha}{2}-1
=2\cdot\left(\frac{\sqrt5}{5}\right)^2-1
=-\frac35
\]
因为 \(0<\alpha<\pi\)，所以 \(\frac{\pi}{2}<\alpha<\pi\)，从而
\(\sin\alpha=\sqrt{1-\cos^2\alpha}=\frac45\).
于是
\[
\sin\left(\alpha-\frac{\pi}{4}\right)
=\sin\alpha\cos\frac{\pi}{4}-\cos\alpha\sin\frac{\pi}{4}
=\frac45\cdot\frac{\sqrt2}{2}-\left(-\frac35\right)\cdot\frac{\sqrt2}{2}
=\frac{7\sqrt2}{10}
\]
故选：D.
\end{solution}
\section{多选题}

\begin{problem}
记 \(S_n\) 为等比数列 \(\{a_n\}\) 的前 \(n\) 项和，\(q\) 为 \(\{a_n\}\) 的公比，\(q>0\). 若 \(S_3=7\)，\(a_3=1\)，则（\quad）

\choices
    {\(q=\frac12\)}
    {\(a_5=\frac19\)}
    {\(S_5=8\)}
    {\(a_n+S_n=8\)}

\end{problem}

\begin{answer}
AD
\end{answer}

\begin{solution}
由题意，\(a_1q^2=1\)，\(a_1+a_1q+a_1q^2=7\).
解得 \(a_1=4\)，\(q=\frac12\)（另一组 \(q<0\) 舍去），故 A 正确.

此时
\[
a_5=a_1q^4=4\cdot\left(\frac12\right)^4=\frac14
\]
故 B 错误.

又
\[
S_5=\frac{a_1(1-q^5)}{1-q}
=\frac{4\left(1-\frac{1}{32}\right)}{1-\frac12}
=\frac{31}{4}
\]
故 C 错误.

并且
\(a_n=4\left(\frac12\right)^{n-1}=2^{3-n}\)，\(S_n=\frac{4\left[1-\left(\frac12\right)^n\right]}{1-\frac12}=8-2^{3-n}\)，
所以
\(a_n+S_n=2^{3-n}+8-2^{3-n}=8\)，
故 D 正确. 故选：AD.
\end{solution}

\begin{problem}
已知 \(f(x)\) 是定义在 \(\R\) 上的奇函数，且当 \(x>0\) 时，\(f(x)=(x^2-3)\e^x+2\)，则（\quad）

\choices
    {\(f(0)=0\)}
    {当 \(x<0\) 时，\(f(x)=-(x^2-3)\e^{-x}-2\)}
    {\(f(x)\ge 2\) 当且仅当 \(x\ge \sqrt3\)}
    {\(x=-1\) 是 \(f(x)\) 的极大值点}

\end{problem}

\begin{answer}
ABD
\end{answer}

\begin{solution}
A：因为 \(f(x)\) 是定义在 \(\R\) 上的奇函数，所以 \(f(0)=0\)，A 正确.

B：当 \(x<0\) 时，\(-x>0\)，故
\[
f(x)=-f(-x)=-\left[((-x)^2-3)\e^{-x}+2\right]=-(x^2-3)\e^{-x}-2
\]
B 正确.

C：由 B，
\[
f(-1)=-(1-3)\e-2=2\e-2=2(\e-1)>2
\]
故 C 错误.

D：当 \(x<0\) 时，
\(f(x)=(3-x^2)\e^{-x}-2\)，
于是
\(f'(x)=(x^2-2x-3)\e^{-x}\).
令 \(f'(x)=0\)，得 \(x=-1\) 或 \(x=3\)（舍去）. 当 \(x\in(-\infty,-1)\) 时，\(f'(x)>0\)；当 \(x\in(-1,0)\) 时，\(f'(x)<0\). 因此 \(x=-1\) 是极大值点，D 正确. 故选：ABD.
\end{solution}

\begin{problem}
双曲线 \(C:\frac{x^2}{a^2}-\frac{y^2}{b^2}=1\)（\(a>0\)，\(b>0\)）的左、右焦点分别是 \(F_1,F_2\)，左、右顶点分别为 \(A_1,A_2\). 以 \(F_1F_2\) 为直径的圆与 \(C\) 的一条渐近线交于 \(M,N\) 两点，且 \(\angle NA_1M=\frac{5\pi}{6}\)，则（\quad）

\choices
    {\(\angle A_1MA_2=\frac{\pi}{6}\)}
    {\(|MA_1|=2|MA_2|\)}
    {\(C\) 的离心率为 \(\sqrt{13}\)}
    {当 \(a=\sqrt2\) 时，四边形 \(NA_1MA_2\) 的面积为 \(8\sqrt3\)}

\end{problem}

\begin{answer}
ACD
\end{answer}

\begin{solution}
不妨设所取渐近线为 \(y=\frac{b}{a}x\)，且 \(M\) 在第一象限，\(N\) 在第三象限.

\bitmapfigure[max width=3.8cm,max totalheight=4.8cm]{img/2025/national_paper_2/q11_solution_fig1.png}

由双曲线的中心对称性，四边形 \(A_1MA_2N\) 是平行四边形，所以
\[
\angle A_1MA_2=\pi-\angle NA_1M=\pi-\frac{5\pi}{6}=\frac{\pi}{6}
\]
故 A 正确.

\bitmapfigure[max width=3.6cm,max totalheight=4.8cm]{img/2025/national_paper_2/q11_solution_fig2.png}

设 \(O\) 为原点. 因为 \(M\) 在以 \(F_1F_2\) 为直径的圆上，所以 \(|OM|=c\). 由渐近线方程和圆的方程可得 \(M\) 的横坐标等于 \(a\)，故 \(|MA_2|=b\). 又在直角三角形 \(A_1A_2M\) 中，由 \(\angle A_1MA_2=\frac{\pi}{6}\) 得
\(2|MA_2|=\sqrt3|MA_1|\)，
所以 B 错误.

由上式和 \(|MA_2|=b\)，再结合 \(|OM|=c\) 可推出
\(c^2=13a^2\)，
故双曲线的离心率
\(e=\frac{c}{a}=\sqrt{13}\)，
C 正确.

当 \(a=\sqrt2\) 时，由 \(e=\sqrt{13}\) 得 \(c=\sqrt{26}\)，再由 \(c^2=a^2+b^2\) 得
\(
b=2\sqrt6
\)
四边形 \(NA_1MA_2\) 为平行四边形，其面积为
\(
2S_{\triangle MA_1A_2}=2\times \frac12 \times 2\sqrt2 \times 2\sqrt6=8\sqrt3
\)
故 D 正确. 故选：ACD.
\end{solution}
\section{填空题}

\begin{problem}
已知平面向量 \(\bs{a}=(x,1)\)，\(\bs{b}=(x-1,2x)\)，若 \(\bs{a}\perp(\bs{a}-\bs{b})\)，则 \(|\bs{a}|=\)\fillinblank{}.

\end{problem}

\begin{answer}
\(\sqrt2\)
\end{answer}

\begin{solution}
由题意，
\(\bs{a}-\bs{b}=(1,1-2x)\).
又 \(\bs{a}\perp(\bs{a}-\bs{b})\)，故
\(\bs{a}\cdot(\bs{a}-\bs{b})=0\)，
即
\(x+1-2x=0\).
解得 \(x=1\). 所以 \(\bs{a}=(1,1)\)，故
\(|\bs{a}|=\sqrt{1^2+1^2}=\sqrt2\).
故答案为：\(\sqrt2\).
\end{solution}

\begin{problem}
若 \(x=2\) 是函数 \(f(x)=(x-1)(x-2)(x-a)\) 的极值点，则 \(f(0)=\)\fillinblank{}.

\end{problem}

\begin{answer}
\(-4\)
\end{answer}

\begin{solution}
因为
\(f'(x)=(x-a)(x-1)+(x-1)(x-2)+(x-a)(x-2)\).
又 \(x=2\) 是极值点，所以 \(f'(2)=0\)，即
\(2-a=0\).
故 \(a=2\). 于是
\(f(0)=(-1)\times(-2)\times(-2)=-4\).
故答案为：\(-4\).
\end{solution}

\begin{problem}
一个底面半径为 \(4\text{cm}\)、高为 \(9\text{cm}\) 的封闭圆柱形容器内有两个半径相等的铁球，则铁球半径的最大值为\fillinblank{}\(\text{cm}\).

\end{problem}

\begin{answer}
\(2.5\)
\end{answer}

\begin{solution}
设铁球半径为 \(r\)，则两球球心距离为 \(2r\). 由题意可得
\((2r)^2=(8-2r)^2+(9-2r)^2\).
化简得
\(4r^2-68r+145=0\)，
即
\((2r-5)(2r-29)=0\).
又 \(r<4\)，故 \(r=2.5\). 故答案为：\(2.5\).

\bitmapfigure[max width=4.2cm,max totalheight=4.8cm]{img/2025/national_paper_2/q14_fig1.png}
\end{solution}
\section{解答题}

\begin{problem}
已知函数 \(f(x)=\cos(2x+\varphi)\)（\(0\le \varphi<\pi\)），且 \(f(0)=\frac12\).
\begin{enumerate}
\item 求 \(\varphi\)；
\item 设函数 \(g(x)=f(x)+f\left(x-\frac{\pi}{6}\right)\)，求 \(g(x)\) 的值域和单调区间.
\end{enumerate}

\end{problem}

\begin{solution}
（1）由 \(f(0)=\cos\varphi=\frac12\)，且 \(0\le \varphi<\pi\)，可得
\(\varphi=\frac{\pi}{3}\).

（2）由（1）得
\(f(x)=\cos\left(2x+\frac{\pi}{3}\right)\).
于是
\(
g(x)=\cos\left(2x+\frac{\pi}{3}\right)+\cos 2x
=\sqrt3\cos\left(2x+\frac{\pi}{6}\right)
\)
所以 \(g(x)\) 的值域为
\([-\sqrt3,\sqrt3]\).

当
\(2k\pi\le 2x+\frac{\pi}{6}\le (2k+1)\pi\)
时，函数 \(g(x)\) 单调递减，解得
\(-\frac{\pi}{12}+k\pi\le x\le \frac{5\pi}{12}+k\pi\)，\(k\in\Z\).

当
\((2k+1)\pi\le 2x+\frac{\pi}{6}\le (2k+2)\pi\)
时，函数 \(g(x)\) 单调递增，解得
\(\frac{5\pi}{12}+k\pi\le x\le \frac{11\pi}{12}+k\pi\)，\(k\in\Z\).
\end{solution}

\begin{problem}
已知椭圆 \(C:\frac{x^2}{a^2}+\frac{y^2}{b^2}=1\)（\(a>b>0\)）的离心率为 \(\frac{\sqrt2}{2}\)，长轴长为 4.
\begin{enumerate}
\item 求 \(C\) 的方程；
\item 过点 \((0,-2)\) 的直线 \(l\) 与 \(C\) 交于 \(A,B\) 两点，\(O\) 为坐标原点. 若 \(\triangle OAB\) 的面积为 \(\sqrt2\)，求 \(|AB|\).
\end{enumerate}

\end{problem}

\begin{solution}
（1）由长轴长为 4，得 \(a=2\). 又离心率为 \(\frac{\sqrt2}{2}\)，所以
\(c=ae=2\cdot \frac{\sqrt2}{2}=\sqrt2\).
故
\(b^2=a^2-c^2=4-2=2\).
所以椭圆方程为
\(\frac{x^2}{4}+\frac{y^2}{2}=1\).

（2）由题设直线 \(AB\) 的斜率不为 0，设直线 \(l\) 的方程为
\(x=t(y+2)\)，

\bitmapfigure[max width=4.0cm,max totalheight=4.8cm]{img/2025/national_paper_2/q16_solution_fig1.png}

其中 \(A(x_1,y_1)\)，\(B(x_2,y_2)\). 联立 \(x=t(y+2)\) 与 \(x^2+2y^2=4\)
可得
\((t^2+2)y^2+4t^2y+4t^2-4=0\).
因此
\(y_1+y_2=-\frac{4t^2}{t^2+2}\)，\(y_1y_2=\frac{4t^2-4}{t^2+2}\).
又
\(
\Delta=16t^4-4(t^2+2)(4t^2-4)=4(8-4t^2)>0
\)
故
\(-\sqrt2<t<\sqrt2\).
又
\[
S_{\triangle OAB}=\frac12\cdot |2t|\cdot |y_1-y_2|
=|t|\sqrt{(y_1+y_2)^2-4y_1y_2}
=\frac{|t|\sqrt{32-16t^2}}{t^2+2}
=\sqrt2
\]
解得
\(t=\pm \frac{\sqrt6}{3}\).
于是
\[
|AB|=\sqrt{1+t^2}\,|y_1-y_2|
=\sqrt{\frac53}\cdot\frac{\sqrt{32-16\times\frac23}}{\frac23+2}
=\sqrt5
\]
\end{solution}

\begin{problem}
如图，在四边形 \(ABCD\) 中，\(AB//CD\)，\(\angle DAB=90^\circ\)，\(F\) 为 \(CD\) 的中点，点 \(E\) 在 \(AB\) 上，\(EF//AD\)，\(AB=3AD\)，\(CD=2AD\).

将四边形 \(EFDA\) 沿 \(EF\) 翻折至四边形 \(EFD'A'\)，使得面 \(EFD'A'\) 与面 \(EFCB\) 所成的二面角为 \(60^\circ\).
\begin{enumerate}
\item 证明：\(A'B//\) 平面 \(CD'F\)；
\item 求面 \(BCD'\) 与面 \(EFD'A'\) 所成的二面角的正弦值.
\end{enumerate}

\bitmapfigure[max width=6.0cm,max totalheight=4.8cm]{img/2025/national_paper_2/q17_fig1.png}

\end{problem}

\begin{solution}
（1）设 \(AD=1\)，则 \(AB=3\)，\(CD=2\). 因为 \(F\) 为 \(CD\) 的中点，所以 \(DF=1\). 又因为 \(EF//AD\)，\(AB//CD\)，所以 \(AEFD\) 是平行四边形，从而 \(AE//DF\). 翻折后仍有 \(A'E//D'F\).

因为 \(D'F\subset\) 平面 \(CD'F\)，且 \(A'E\not\subset\) 平面 \(CD'F\)，故 \(A'E//\) 平面 \(CD'F\). 又因为 \(FC//EB\)，且 \(FC\subset\) 平面 \(CD'F\)，\(EB\not\subset\) 平面 \(CD'F\)，故 \(EB//\) 平面 \(CD'F\).

又 \(A'E\cap EB=E\)，且 \(A'E,EB\subset\) 平面 \(A'EB\)，所以平面 \(A'EB//\) 平面 \(CD'F\). 因此
\(A'B//\text{平面 }CD'F\).

（2）以 \(F\) 为原点，\(FE,FC\) 以及垂直于平面 \(EFCB\) 的直线分别为 \(x,y,z\) 轴，建立空间直角坐标系.

\bitmapfigure[max width=4.0cm,max totalheight=4.8cm]{img/2025/national_paper_2/q17_solution_fig1.png}

由题意可得
\(B(1,2,0)\)，\(C(0,1,0)\)，\(D'\left(0,\frac12,\frac{\sqrt3}{2}\right)\)，\(E(1,0,0)\)，\(F(0,0,0)\).
于是 \(\overrightarrow{BC}=(-1,-1,0)\)，\(\overrightarrow{CD'}=\left(0,-\frac12,\frac{\sqrt3}{2}\right)\)，\(\overrightarrow{FE}=(1,0,0)\)，\(\overrightarrow{FD'}=\left(0,\frac12,\frac{\sqrt3}{2}\right)\).

设平面 \(BCD'\) 的法向量为 \(\bs{n}=(x,y,z)\)，则 \(\overrightarrow{BC}\cdot \bs{n}=0\)，\(\overrightarrow{CD'}\cdot \bs{n}=0\).
可取
\(\bs{n}=(-\sqrt3,\sqrt3,1)\).

设平面 \(EFD'A'\) 的法向量为 \(\bs{m}=(x_1,y_1,z_1)\)，则 \(\overrightarrow{FE}\cdot \bs{m}=0\)，\(\overrightarrow{FD'}\cdot \bs{m}=0\).
可取
\(\bs{m}=(0,\sqrt3,-1)\).

所以两平面的夹角余弦为
\(
\cos\theta=\frac{|\bs{m}\cdot \bs{n}|}{|\bs{m}|\,|\bs{n}|}=\frac{1}{\sqrt7}
\)
故所求二面角的正弦值为
\(
\sqrt{1-\frac17}=\frac{\sqrt{42}}{7}
\)
\end{solution}

\begin{problem}
已知函数 \(f(x)=\ln(1+x)-x+\frac12x^2-kx^3\)，其中 \(0<k<\frac13\).
\begin{enumerate}
\item 证明：\(f(x)\) 在区间 \((0,+\infty)\) 存在唯一的极值点和唯一的零点；
\item 设 \(x_1,x_2\) 分别为 \(f(x)\) 在区间 \((0,+\infty)\) 的极值点和零点.
    \begin{enumerate}
    \item 设函数 \(g(t)=f(x_1+t)-f(x_1-t)\)，证明：\(g(t)\) 在区间 \((0,x_1)\) 单调递减；
    \item 比较 \(2x_1\) 与 \(x_2\) 的大小，并证明你的结论.
    \end{enumerate}
\end{enumerate}

\end{problem}

\begin{solution}
（1）由题意，
\(
f'(x)=\frac{1}{1+x}-1+x-3kx^2=x^2\left(\frac{1}{1+x}-3k\right)
\)
设
\(h(x)=\frac{1}{1+x}-3k\)，\(x>0\)，
则
\(h'(x)=-\frac{1}{(1+x)^2}<0\).
所以 \(h(x)\) 在 \((0,+\infty)\) 上单调递减. 又 \(h(0)=1-3k>0\)，且存在唯一
\(x_1=\frac{1}{3k}-1\)
使 \(h(x_1)=0\). 故 \(f'(x)\) 在 \((0,x_1)\) 上为正，在 \((x_1,+\infty)\) 上为负，所以 \(f(x)\) 在 \((0,+\infty)\) 上存在唯一极值点.

再看函数
\(\ln(1+x)-x\)，
其导数为
\(\frac{1}{1+x}-1=-\frac{x}{1+x}<0\)，
故 \(\ln(1+x)-x<0\)（\(x>0\)）. 又当 \(x\to +\infty\) 时，\(\frac12x^2-kx^3\to -\infty\)，所以
\(f(x)\to -\infty\).
由 \(h(x)=\frac{1}{1+x}-3k\) 且 \(h(x_1)=0\) 得，当 \(0<x<x_1\) 时，
\(\frac{1}{1+x}-3k>\frac{1}{1+x_1}-3k=0\)，
所以 \(f'(x)>0\)；当 \(x>x_1\) 时，
\(\frac{1}{1+x}-3k<\frac{1}{1+x_1}-3k=0\)，
所以 \(f'(x)<0\). 因此 \(f(x)\) 在 \((0,x_1)\) 上递增，在 \((x_1,+\infty)\) 上递减.

又 \(f(0)=0\)，并且对充分大的 \(x\)，有 \(f(x)<0\). 由于 \(f(x)\) 先增后减，故其图象在 \((0,+\infty)\) 上与 \(x\) 轴恰有一个交点，所以 \(f(x)\) 在 \((0,+\infty)\) 上恰有一个零点，记为 \(x_2\).

（2）（1）由（1）知
\(x_1+1=\frac{1}{3k}\)，\(f'(x)=x^2\left(\frac{1}{1+x}-\frac{1}{1+x_1}\right)\)
因为
\(g(t)=f(x_1+t)-f(x_1-t)\)，
所以
\(g'(t)=f'(x_1+t)-f'(x_1-t)\).
代入得
\(
g'(t)=(x_1+t)^2\left(\frac{1}{x_1+t+1}-\frac{1}{1+x_1}\right)
+(x_1-t)^2\left(\frac{1}{x_1-t+1}-\frac{1}{1+x_1}\right)
\)
所以
\(
g'(t)=-\frac{t(x_1+t)^2}{(x_1+t+1)(x_1+1)}
+\frac{t(x_1-t)^2}{(x_1-t+1)(x_1+1)}
\)
又由 \(x_1+1=\frac{1}{3k}\) 可得
\(
g'(t)=3kt\left[-\frac{(x_1+t)^2}{x_1+t+1}+\frac{(x_1-t)^2}{x_1-t+1}\right]
=\frac{6kt^2(t^2-x_1^2-2x_1)}{(1+x_1)^2-t^2}
\)
当 \(t\in(0,x_1)\) 时，\(k>0\)，\(t^2-x_1^2-2x_1<0\)，\((1+x_1)^2-t^2>0\)，故
\(g'(t)<0\).
因此 \(g(t)\) 在 \((0,x_1)\) 上单调递减.

（2）由（1）知 \(g(t)\) 在 \((0,x_1)\) 上单调递减，所以
\(g(0)>g(x_1)\).
即
\(0>f(2x_1)\).
又 \(f(x_2)=0\)，并且由（1）知 \(f(x)\) 在 \((x_1,+\infty)\) 上单调递减，因此必有
\(2x_1>x_2\).
\end{solution}

\begin{problem}
甲、乙两人进行乒乓球练习，每个球胜者得 1 分，负者得 0 分. 设每个球甲胜的概率为 \(p\)（\(\frac12<p<1\)），乙胜的概率为 \(q\)，\(p+q=1\)，且各球的胜负相互独立. 对正整数 \(k\ge 2\)，记 \(p_k\) 为打完 \(k\) 个球后甲比乙至少多得 2 分的概率，\(q_k\) 为打完 \(k\) 个球后乙比甲至少多得 2 分的概率.
\begin{enumerate}
\item 求 \(p_3,p_4\)（用 \(p\) 表示）；
\item 若 \(\frac{p_4-p_3}{q_4-q_3}=4\)，求 \(p\)；
\item 证明：对任意正整数 \(m\)，\(p_{2m+1}-q_{2m+1}<p_{2m}-q_{2m}<p_{2m+2}-q_{2m+2}\).
\end{enumerate}

\end{problem}

\begin{solution}
（1）打完 3 个球后甲比乙至少多得 2 分，只能是甲胜 3 场，所以
\(p_3=p^3\).
打完 4 个球后甲比乙至少多得 2 分，只能是甲胜 3 场或 4 场，所以
\[
p_4=\binom43 p^3q+\binom44 p^4
=4p^3(1-p)+p^4
=p^3(4-3p)
\]

（2）打完 3 个球后乙比甲至少多得 2 分，只能是乙胜 3 场，所以
\(q_3=q^3\).
打完 4 个球后乙比甲至少多得 2 分，只能是乙胜 3 场或 4 场，所以
\[
q_4=\binom43 q^3p+\binom44 q^4
=4q^3(1-q)+q^4
=q^3(4-3q)
\]
由
\(\frac{p_4-p_3}{q_4-q_3}=4\)
得
\(\frac{3p^3(1-p)}{3q^3(1-q)}=4\).
又 \(p+q=1\)，所以
\(\left(\frac{p}{q}\right)^2=4\).
因为 \(p,q>0\)，故 \(p=2q=2(1-p)\)，解得
\(p=\frac23\).

（3）我们有
\[
p_{2m}-p_{2m+1}
=\sum_{k=0}^{m-1}\binom{2m}{k}p^{2m-k}q^k
-\sum_{k=0}^{m-1}\binom{2m+1}{k}p^{2m+1-k}q^k
\]
因为
\(\binom{2m+1}{k}=\binom{2m}{k}+\binom{2m}{k-1}\)，
所以
\begin{align*}
p_{2m}-p_{2m+1}
&=\sum_{k=0}^{m-1}\binom{2m}{k}p^{2m-k}q^k
-\sum_{k=0}^{m-1}\binom{2m}{k}p^{2m+1-k}q^k
-\sum_{k=0}^{m-1}\binom{2m}{k-1}p^{2m+1-k}q^k \\
&=(1-p)\sum_{k=0}^{m-1}\binom{2m}{k}p^{2m-k}q^k
-\sum_{k=0}^{m-1}\binom{2m}{k-1}p^{2m+1-k}q^k \\
&=\sum_{k=0}^{m-1}\binom{2m}{k}p^{2m-k}q^{k+1}
-\sum_{k=0}^{m-2}\binom{2m}{k}p^{2m-k}q^{k+1} \\
&=\binom{2m}{m-1}p^{m+1}q^m
\end{align*}
同理可得
\(p_{2m+2}-p_{2m+1}=\binom{2m+1}{m}p^{m+2}q^m\)，
\(q_{2m}-q_{2m+1}=\binom{2m}{m-1}q^{m+1}p^m\)，
\(q_{2m+2}-q_{2m+1}=\binom{2m+1}{m}q^{m+2}p^m\).
于是
\(
p_{2m}-p_{2m+1}
=p\left(\binom{2m}{m-1}p^mq^m\right)
>q\left(\binom{2m}{m-1}q^mp^m\right)
=q_{2m}-q_{2m+1}
\)
从而
\(p_{2m+1}-q_{2m+1}<p_{2m}-q_{2m}\).

另一方面，
\begin{align*}
p_{2m+2}-p_{2m}
&=(p_{2m+2}-p_{2m+1})-(p_{2m}-p_{2m+1}) \\
&=\binom{2m+1}{m}p^{m+2}q^m-\binom{2m}{m-1}p^{m+1}q^m \\
&=p^mq^m\cdot p\left(p\binom{2m+1}{m}-\binom{2m}{m-1}\right) \\
&=\frac{(2m+1)!}{m!(m+1)!}p^mq^m\cdot p\left(p-\frac{m}{2m+1}\right)
\end{align*}
同理
\(q_{2m+2}-q_{2m} =\frac{(2m+1)!}{m!(m+1)!}p^mq^m\cdot q\left(q-\frac{m}{2m+1}\right)\).
作差得
\begin{align*}
&(p_{2m+2}-p_{2m})-(q_{2m+2}-q_{2m}) \\
&\quad=\frac{(2m+1)!}{m!(m+1)!}p^mq^m
\left[p\left(p-\frac{m}{2m+1}\right)-q\left(q-\frac{m}{2m+1}\right)\right] \\
&\quad=\frac{(2m+1)!}{m!(m+1)!}p^mq^m\cdot \frac{m+1}{2m+1}(p-q)>0
\end{align*}
因此
\(p_{2m}-q_{2m}<p_{2m+2}-q_{2m+2}\).
综上，
\(
p_{2m+1}-q_{2m+1}<p_{2m}-q_{2m}<p_{2m+2}-q_{2m+2}
\)
证毕.
\end{solution}
